\chapter{Le métier des produits structurés}\label{ch:dv:the-structured-products-business}

Un client de banque privée achète une note à deux ans sur un indice. Elle lui rend l'intégralité de son capital à l'échéance, plus 70~\% de la
hausse de l'indice. L'option qui verse les 70~\% est une option d'achat ordinaire à la monnaie, et, à la volatilité de l'indice et au
taux courant, elle coûte environ 12.5 pour 100. Ce qui déterminait si 70 était possible n'était pas cette option. C'était le prix auquel
la banque qui a émis la note emprunte. Une note est une obligation de son émetteur, et un émetteur qui emprunte à un spread élevé
au-dessus du taux sans risque actualise plus fortement les 100 garantis, ce qui laisse plus d'argent du client pour acheter des options.
Sans spread, la note aurait pu offrir 35~\%. À un spread de 2.4~\% par an, elle peut offrir 70~\%. Le client a été rémunéré pour
prêter à la banque, et peut-être sans le savoir. Ce chapitre démonte le métier~: ses produits et ses enveloppes, la
décomposition qui évalue chaque note, la distribution et ses marges et règles, le financement de l'émetteur, et les
indices systématiques sur lesquels les produits sont de plus en plus adossés.

\section{Produits et enveloppes}

\begin{definition}[Produit structuré]\label{def:dv:the-structured-products-business:product}
Un \emph{produit structuré}\index{produit structuré} est un placement dont le gain est une fonction prédéfinie d'un ou de plusieurs
sous-jacents, construit en combinant un titre à revenu fixe avec des dérivés et vendu comme un seul ensemble.
\end{definition}

\begin{definition}[Note structurée]\label{def:dv:the-structured-products-business:note}
Une \emph{note structurée}\index{note structurée} est un \omterm{def:dv:the-structured-products-business:product}{produit structuré} émis sous forme de titre de dette d'une banque ou de son
véhicule d'émission~: l'investisseur porte le risque de crédit de l'émetteur autant que le risque de marché du gain.
\end{definition}

Le même gain peut être vendu dans plusieurs enveloppes. La note est la plus courante. Un certificat ou un warrant coté sur une bourse
reste une dette de son émetteur. Le gain peut aussi se loger dans un fonds, le collatéral étant conservé séparément, ou dans un
dépôt, souvent assorti de garanties sur une partie du capital. Il peut être écrit comme un swap de gré à gré pour des institutions, dont
le risque de contrepartie est collatéralisé. Il peut même être emballé sous forme de fonds indiciel coté (ETF) (livre~1, chapitre~14) ou un
total return swap (livre~1, chapitre~6) sur un indice de stratégie. L'enveloppe détermine qui porte le risque de crédit de l'émetteur, comment
le produit est taxé et réglementé, et si un marché secondaire existe. Le gain se décompose de la même façon dans tous les cas.

Les deux archétypes sont la \omterm{def:dv:the-structured-products-business:protected}{note à capital protégé} et le produit à rendement amélioré.

\begin{definition}[Note à capital protégé]\label{def:dv:the-structured-products-business:protected}
Une \emph{note à capital protégé}\index{note à capital protégé} rembourse au minimum une part annoncée de son notionnel à l'échéance,
par exemple 90~\% ou 100~\%, quelle que soit l'évolution du sous-jacent, plus une part de la hausse du sous-jacent~; la protection ne vaut que
ce que vaut le crédit de l'émetteur.
\end{definition}

\begin{definition}[Reverse convertible]\label{def:dv:the-structured-products-business:rc}
Un \emph{reverse convertible}\index{reverse convertible} verse un coupon supérieur au taux de financement de l'émetteur et rembourse le pair à
l'échéance, sauf si le sous-jacent est passé sous le strike, auquel cas il rembourse la valeur du sous-jacent (ou des actions)~:
l'investisseur a vendu un put.
\end{definition}

\begin{definition}[Taux de participation]\label{def:dv:the-structured-products-business:participation}
Le \emph{taux de participation}\index{taux de participation} d'une note est la part de la performance du sous-jacent qu'elle verse au-delà
du montant protégé, $p$ dans $100\,(\text{protection})+100\,p\,(S_T/S_0-1)^+$.
\end{definition}

Entre les archétypes se placent les \omterm{def:dv:autocallables:autocallable}{autocallables} du chapitre 18, qui échangent la protection du capital contre des coupons, et une longue
série de variantes. Il existe des caps et des floors sur la participation, un moyennage du niveau final (chapitre 16), des worst-of sur paniers
(chapitre 17), et des barrières knock-in (chapitre 15). Chaque variante modifie la jambe optionnelle, et chacune a un prix que
la décomposition rend visible.

\section{Anatomie d'une note~: une obligation plus des options}

Toute note est évaluée par la même identité. Le prix d'émission, 100, est égal à l'obligation zéro-coupon de l'émetteur pour les
montants garantis, plus les options, plus la marge~:
\[
100=\underbrace{\pi\cdot100\,e^{-(r+s)T}}_{\text{obligation au taux de l'émetteur}}+\underbrace{p\cdot C}_{\text{options}}+
\underbrace{m}_{\text{marge}},
\]
avec $\pi$ la part protégée, $r$ le taux sans risque, $s$ le spread de financement de l'émetteur (livre~1, chapitre~17), $C$ le prix de
la jambe optionnelle par unité de participation, et $m$ la \omterm{def:dv:the-structured-products-business:margin}{marge de structuration}. Étant donnés trois quelconques parmi $p$, $s$, $m$ et les caractéristiques
de l'option, l'identité donne la quatrième.

\begin{definition}[Marge de structuration]\label{def:dv:the-structured-products-business:margin}
La \emph{marge de structuration}\index{marge de structuration} est l'écart entre le prix d'émission d'une note et la valeur de ses
composantes évaluées sur la courbe de financement propre de l'émetteur et aux prix d'options du desk~: le revenu partagé entre l'émetteur et les
distributeurs.
\end{definition}

\begin{example}[Participation de soixante-dix pour cent]\label{ex:dv:the-structured-products-business:participation}
Une note à deux ans, protégée à 100~\%, sur un indice avec un taux de 3~\% et un rendement de dividende de 1.5~\%, avec la volatilité à deux ans
à la monnaie du chapitre 9, soit 20.5~\%. L'option d'achat à la monnaie coûte 12.49 pour 100, et la marge est de 1.5. À spread de financement nul
l'obligation coûte 94.18, ce qui laisse 4.32 pour les options~: une participation de 34.6~\%. Chaque demi-point de spread ajoute environ 7.5
points de participation. À 1~\%, elle est de 49.6~\%, à 2~\%, de 64.2~\%, et à 2.40~\%, l'obligation coûte 89.76 et la participation atteint
70~\% (\cref{fig:dv:the-structured-products-business:participation}).
\end{example}

\begin{omfigure}
\begin{tikzpicture}
\begin{axis}[width=10.5cm,height=5cm,
  xlabel={spread de financement de l'émetteur (points de base par an)},ylabel={participation (\%)},
  tick label style={font=\small},label style={font=\small},
  xmin=0,xmax=300,ymin=30,ymax=82,grid=major,grid style={gray!25},table/col sep=comma]
\addplot[omThm,very thick,mark=*,mark size=1.4pt] table[x=spread,y=participation]{figdata/derivatives/19-the-structured-products-business/participation.csv};
\draw[gray,dashed] (axis cs:0,70) -- (axis cs:240.3,70) -- (axis cs:240.3,30);
\node[font=\footnotesize,anchor=south east] at (axis cs:235,70.5) {70~\% à 240~pb};
\end{axis}
\end{tikzpicture}

\omcaption{La participation qu'une note à deux ans protégée à 100~\% peut offrir, après une marge de 1.5~\%, selon le
spread de financement de l'émetteur (taux de 3~\%, rendement de dividende de 1.5~\%, volatilité à deux ans du chapitre 9). Le spread, et non l'option, fixe la participation~:
sans lui la note pourrait offrir 35~\%. Données~: le tutoriel.}\label{fig:dv:the-structured-products-business:participation}
\end{omfigure}

Le spread est payé par l'investisseur. Un spread plus élevé signifie un émetteur plus fragile, dont le défaut emporterait la protection avec lui. Le
70~\% rémunère ce risque, et une note d'un émetteur plus solide offrant 50~\% n'est pas une moins bonne affaire. La même identité
évalue les autres formes. Avec un spread de 1~\%, le budget achète une participation de 100~\% jusqu'à un cap de 113.3~\%, soit un rendement à deux ans plafonné à 13.3~\%,
au lieu de 49.6~\% sans cap (\cref{fig:dv:the-structured-products-business:payoffs}). Sans spread, le cap descend à 108.7~\%. Le
\omterm{def:dv:the-structured-products-business:rc}{reverse convertible} inverse le sens des flux~: l'investisseur vend un put, et la prime s'ajoute au coupon.

\begin{example}[Un reverse convertible à un an]\label{ex:dv:the-structured-products-business:rc}
Strike 90~\%, un an, spread de l'émetteur de 1~\%, même marché. L'investisseur vend $100/0.9$ puts de strike 90~\%, valant 4.26, et le
gain de financement de l'émetteur vaut 3.92. Après une marge de 1.5, la note verse un coupon de 6.96~\%, contre un taux sans risque de 3~\%. Si
l'indice termine sous 90~\%, l'investisseur reçoit $100\,S_T/(0.9\,S_0)$ au lieu du pair, et le coupon.
\end{example}

\begin{omfigure}
\begin{tikzpicture}
\begin{axis}[width=10.5cm,height=5cm,
  xlabel={indice à l'échéance (\% du départ)},ylabel={remboursement (pour 100)},
  tick label style={font=\small},label style={font=\small},
  xmin=50,xmax=160,ymin=90,ymax=145,grid=major,grid style={gray!25},
  legend columns=2,legend cell align=left,legend style={font=\footnotesize,at={(0.5,-0.3)},anchor=north,draw=none,fill=none,/tikz/every even column/.append style={column sep=8pt}},
  table/col sep=comma]
\addplot[omThm,very thick] table[x=perf,y=protected]{figdata/derivatives/19-the-structured-products-business/payoffs.csv};
\addplot[omDef,very thick,dashed] table[x=perf,y=capped]{figdata/derivatives/19-the-structured-products-business/payoffs.csv};
\legend{70~\% de participation (spread de financement de 2.4~\%),100~\% plafonné (spread de 1~\%)}
\end{axis}
\end{tikzpicture}

\omcaption{Deux notes protégées à 100~\% sur deux ans, dont chacune coûte 100 au client~: 70~\% de participation sans cap, ce qui exige un
spread de financement de 2.4~\%, et une pleine participation plafonnée à 113.3~\%, que finance un spread de 1~\%. Données~: le tutoriel.}\label{fig:dv:the-structured-products-business:payoffs}
\end{omfigure}

\section{Distribution, marges et régulation}

Une note passe par plusieurs mains avant d'atteindre son acheteur. Le structureur la conçoit, et le desk de trading la valorise et la couvre.
L'émetteur, souvent une entité juridique distincte, apporte le bilan. Le distributeur, une banque privée, une banque de détail ou un
conseiller indépendant, le vend. La marge de l'identité ci-dessus est partagée entre eux. Le desk conserve une réserve pour les frais de couverture, et
l'émetteur et le distributeur sont rémunérés par le prix de la note.

Les régulateurs se concentrent sur ce que l'acheteur de détail peut voir. Le règlement de l'Union européenne sur les produits d'investissement de détail packagés exige un
court document d'informations clés standardisé pour chaque produit de ce type. L'organisme d'autorégulation américain FINRA a publié des recommandations sur les produits
présentant des «~caractéristiques nouvelles, compliquées ou enchevêtrées, de type dérivé~», parmi lesquels figurent les \omterm{def:dv:the-structured-products-business:note}{notes structurées}.

\begin{dated}{2026-09}{Règles de publication pour les produits structurés de détail}\label{dat:dv:the-structured-products-business:rules}
\emph{Union européenne.} La Regulation (EU) No 1286/2014 sur les produits d'investissement packagés de détail et fondés sur l'assurance (PRIIPs) impose
au fabricant, avant la mise à disposition d'un produit auprès des investisseurs de détail, d'établir un document d'informations clés et de le publier sur
son site. Le document constitue une information précontractuelle~; il doit être exact, équitable, clair et non trompeur, être autonome et distinct
des documents promotionnels, et ne dépasse pas trois pages au format A4. \emph{États-Unis.} FINRA Regulatory Notice 12-03 (janvier
2012), «~Heightened Supervision of Complex Products~», demande aux broker-dealers d'appliquer une surveillance renforcée aux produits tels que
les \omterm{def:dv:the-structured-products-business:note}{notes structurées}, et fait suite à des avis antérieurs sur les \omterm{def:dv:the-structured-products-business:product}{produits structurés}, les \omterm{def:dv:the-structured-products-business:protected}{notes à capital protégé} et les \omterm{def:dv:the-structured-products-business:rc}{reverse convertibles}. Les deux textes
ont été modifiés ou complétés depuis~; il faut consulter les versions en vigueur.
\end{dated}

Pour un quant, le texte réglementaire compte à trois endroits. Les scénarios de performance et les chiffres de coûts de l'information publiée sont calculés à partir
de modèles, et le modèle doit être celui qu'utilise le desk. La marge doit pouvoir se défendre comme fair value. Et la valeur du produit après émission,
le prix auquel l'émetteur rachète les notes, doit découler de la même valorisation que le prix d'émission, diminué d'une fourchette acheteur-vendeur dont
l'investisseur a été informé.

\section{Financement de l'émetteur et cycle de vie du produit}

Une \omterm{def:dv:the-structured-products-business:note}{note structurée} est un financement pour son émetteur. La banque reçoit 100 aujourd'hui et doit verser le paiement final plus tard. Si elle couvre ce paiement avec
des dérivés, ce qui reste est un prêt à son propre spread de financement. Les notes sont ainsi l'une des sources de financement d'une banque, ce qui explique pourquoi
les émetteurs se disputent le taux auquel ils valorisent leurs propres notes. Un émetteur qui valorise sa propre obligation à un spread élevé peut offrir de meilleures
conditions. La note reste équitablement valorisée si le spread correspond à son coût de financement réel.

Après l'émission, la note vit dans les livres du desk comme tout dérivé. Le desk gère les \omterm{def:dv:greeks-and-the-hedging-pnl:greeks}{grecques} de la jambe optionnelle et finance ou non la
jambe obligataire avec la trésorerie, et il cote un prix secondaire, souvent sous la forme du rachat par l'émetteur. Ce prix évolue avec le marché et avec le
spread de l'émetteur~: un élargissement du spread fait baisser la valeur de chaque note en circulation, et un émetteur qui porte ses notes à la fair value
constate un gain sur sa propre dette. Les rachats anticipés, les \omterm{def:dv:autocallables:autocallable}{autocallables} (chapitre~18) et les knock-in modifient le financement détenu par l'émetteur, ce que la trésorerie doit anticiper.

\section{Indices de stratégie systématique et options sur ces indices}

Les notes peuvent aussi être indexées sur un indice fondé sur des règles plutôt que sur un simple indice de marché. Deux constructions de ce type sont conçues pour rendre les options
moins chères.

\begin{definition}[Indice de stratégie systématique]\label{def:dv:the-structured-products-business:ssi}
Un \emph{indice de stratégie systématique}\index{indice de stratégie systématique} est un indice dont le niveau suit une règle publiée de négociation d'un ou
plusieurs sous-jacents, rééquilibré mécaniquement, et sur lequel sont écrites des notes et des options.
\end{definition}

\begin{definition}[Indice à volatilité cible]\label{def:dv:the-structured-products-business:vt}
Un \emph{indice à volatilité cible}\index{indice à volatilité cible} détient une exposition à un sous-jacent égale à une volatilité cible divisée
par une estimation récente de la volatilité du sous-jacent, plafonnée à un levier maximal, le reste en liquidités~; sa volatilité réalisée reste
proche de la cible quel que soit le comportement du sous-jacent.
\end{definition}

\begin{definition}[Indice à décrément]\label{def:dv:the-structured-products-business:decrement}
Un \emph{indice à décrément}\index{indice à décrément} est construit sur un indice de rendement total en déduisant un montant fixe, un pourcentage annuel ou un
nombre de points, à la place des dividendes effectivement versés~: son forward dépend de cette déduction fixe, et non de prévisions
incertaines de dividendes.
\end{definition}

La volatilité cible rend la volatilité de l'indice presque déterministe. Une option sur celui-ci est presque une option de Black--Scholes à la cible,
bon marché lorsque la cible est inférieure à la volatilité du marché. La couverture de l'émetteur ne comporte alors presque aucun risque de volatilité, et sa position en \omterm{def:dv:greeks-and-the-hedging-pnl:greeks}{vega}
ne dépend plus de la volatilité implicite de l'indice. L'\omterm{def:dv:the-structured-products-business:decrement}{indice à décrément} supprime le \omterm{def:dv:dividends-borrow-and-forwards:risk}{risque de dividende} (chapitre~5). L'émetteur connaît
exactement le forward et peut vendre des puts de longue échéance sur lui sans détenir d'exposition aux dividendes. Lorsque le décrément dépasse les
dividendes attendus, l'indice dérive à la baisse par rapport au marché, et les puts sur lui sont plus chers.

\begin{example}[Options sur un indice à volatilité cible de 10~\%]\label{ex:dv:the-structured-products-business:vt}
Sur le \omterm{def:dv:stochastic-volatility:heston}{modèle de Heston} du chapitre~10, calibré sur la surface du chapitre~9, on construit un \omterm{def:dv:the-structured-products-business:vt}{indice à volatilité cible} de 10~\% avec des estimations de volatilité
pondérées exponentiellement et un levier plafonné à 1.5, avec des taux nuls. Sur un an, l'indice brut réalise 19.8~\% en moyenne, avec un écart-type
de 8.2 points selon les trajectoires. L'indice cible réalise 10.1~\%, avec un écart-type de 0.6 point. Un call d'un an à la monnaie coûte
7.54 sur l'indice brut et 4.07 sur l'indice cible, et les volatilités implicites valent 19.0~\% et 10.2~\%. Le skew de l'indice brut, de 21.6~\% en
90 à 16.8~\% en 110, s'estompe presque sur l'indice cible, de 10.6~\% à 9.9~\% (\cref{fig:dv:the-structured-products-business:vt}).
\end{example}

\begin{omfigure}
\begin{tikzpicture}
\begin{groupplot}[group style={group size=2 by 1,horizontal sep=1.6cm},
  width=6.6cm,height=5cm,
  tick label style={font=\small},label style={font=\small},grid=major,grid style={gray!25},table/col sep=comma]
\nextgroupplot[xlabel={jour de bourse},ylabel={niveau~; exposition (\%)},xmin=0,xmax=252,
  title={\small une trajectoire},
  legend columns=1,legend cell align=left,legend style={font=\footnotesize,at={(0.5,-0.32)},anchor=north,draw=none,fill=none}]
\addplot[gray,thick] table[x=day,y=raw]{figdata/derivatives/19-the-structured-products-business/vt_path.csv};
\addplot[omThm,very thick] table[x=day,y=vt]{figdata/derivatives/19-the-structured-products-business/vt_path.csv};
\addplot[omDef,thick,dashed] table[x=day,y=exposure]{figdata/derivatives/19-the-structured-products-business/vt_path.csv};
\legend{indice,indice à volatilité cible de 10~\%,exposition}
\nextgroupplot[xlabel={prix d'exercice},ylabel={volatilité implicite à 1 an (\%)},xmin=87,xmax=113,ymin=5,ymax=25,xtick={90,100,110},
  title={\small calls sur chaque indice},
  legend columns=1,legend cell align=left,legend style={font=\footnotesize,at={(0.5,-0.32)},anchor=north,draw=none,fill=none}]
\addplot[gray,very thick,mark=*,mark size=1.6pt] table[x=k,y=raw]{figdata/derivatives/19-the-structured-products-business/vt_iv.csv};
\addplot[omThm,very thick,mark=square*,mark size=1.6pt] table[x=k,y=vt]{figdata/derivatives/19-the-structured-products-business/vt_iv.csv};
\legend{indice,indice à volatilité cible de 10~\%}
\end{groupplot}
\end{tikzpicture}

\omcaption{Un \omterm{def:dv:the-structured-products-business:vt}{indice à volatilité cible} de 10~\% sur l'indice de Heston du chapitre 10. À gauche~: une trajectoire de l'indice, de l'indice cible et
de l'exposition, qui baisse lorsque la volatilité réalisée monte. À droite~: les volatilités implicites des calls à un an sur chacun~: celles de l'indice cible
sont proches de 10~\% à chaque prix d'exercice. Données~: le tutoriel.}\label{fig:dv:the-structured-products-business:vt}
\end{omfigure}

\begin{example}[Un indice à décrément face aux dividendes]\label{ex:dv:the-structured-products-business:decrement}
Prenons un taux de 3~\%. Un indice de prix versant 3~\% de dividendes a un forward à cinq ans de 100. Un \omterm{def:dv:the-structured-products-business:decrement}{indice à décrément} de 5~\% construit sur le même
indice de rendement total a un forward de 90.5, car il perd 2~\% par an par rapport à celui-ci
(\cref{fig:dv:the-structured-products-business:decrement}). Une option de vente à la monnaie à cinq ans, de volatilité 20~\%, coûte 15.2 pour 100 sur
l'indice de prix et 19.0 sur l'\omterm{def:dv:the-structured-products-business:decrement}{indice à décrément}, soit 24~\% de plus. C'est ce que reçoit un investisseur qui vend le put au sein d'un \omterm{def:dv:autocallables:autocallable}{autocallable} en échange de
la prise en charge de l'écart entre le décrément et les dividendes effectivement versés.
\end{example}

\begin{omfigure}
\begin{tikzpicture}
\begin{axis}[width=10.5cm,height=4.6cm,
  xlabel={années},ylabel={forward (\% d'aujourd'hui)},
  tick label style={font=\small},label style={font=\small},
  xmin=0,xmax=5,ymin=88,ymax=102,grid=major,grid style={gray!25},
  legend columns=2,legend cell align=left,legend style={font=\footnotesize,at={(0.5,-0.32)},anchor=north,draw=none,fill=none,/tikz/every even column/.append style={column sep=8pt}},
  table/col sep=comma]
\addplot[gray,very thick] table[x=t,y=price_fwd]{figdata/derivatives/19-the-structured-products-business/decrement.csv};
\addplot[omThm,very thick] table[x=t,y=dec_fwd]{figdata/derivatives/19-the-structured-products-business/decrement.csv};
\legend{indice de prix (3~\% de dividendes),indice à décrément de 5~\%}
\end{axis}
\end{tikzpicture}

\omcaption{Forwards d'un indice de prix versant 3~\% de dividendes et d'un \omterm{def:dv:the-structured-products-business:decrement}{indice à décrément} de 5~\% construit sur le même indice de rendement total, avec un taux de 3~\%.
Le décrément fixe le forward, et l'indice dérive à la baisse par rapport au marché d'un montant égal à l'excédent du décrément sur les dividendes. Données~:
le code du chapitre.}\label{fig:dv:the-structured-products-business:decrement}
\end{omfigure}

Ces deux constructions transfèrent un risque de l'émetteur vers l'investisseur. L'\omterm{def:dv:the-structured-products-business:vt}{indice à volatilité cible} fait supporter à l'investisseur le risque de calendrier de son rééquilibrage,
qui achète après les hausses et vend après les baisses, une configuration de gamma court. Le décrément fait supporter à l'investisseur l'écart entre le décrément et
les dividendes. Le profil de gains semble le même que sur l'indice nu, et il ne l'est pas. La première question d'un structureur sur un indice de stratégie est
donc ce qu'il a retiré du prix de l'option, et qui le détient désormais.

\section{Tutoriel~: décomposer une note}\label{tut:dv:the-structured-products-business:tut}

\begin{tutorial}
\textbf{Objectif.} Décomposer une note protégée en une obligation au taux de l'émetteur et une jambe d'option~; résoudre le \omterm{def:dv:the-structured-products-business:participation}{taux de participation}, le plafond, le
coupon du \omterm{def:dv:the-structured-products-business:rc}{reverse convertible} et le spread de financement qu'un objectif exige~; construire un \omterm{def:dv:the-structured-products-business:vt}{indice à volatilité cible} et un \omterm{def:dv:the-structured-products-business:decrement}{indice à décrément}, et pricer des
options sur ceux-ci. \textbf{État final~:} les quatre figures et les nombres du problème du week-end.
\begin{tutsteps}
\item \textbf{L'identité budgétaire, dans les deux sens}~:
\omcode{code/firm/termsheet/firm_termsheet.py}{57}{72}{Taux de participation à partir du spread, et spread à partir d'un taux de participation.}\label{lst:dv:the-structured-products-business:budget}
\item \textbf{L'\omterm{def:dv:the-structured-products-business:vt}{indice à volatilité cible}}, avec une estimation de volatilité retardée pour que la règle n'utilise que des données passées~:
\omcode{code/firm/termsheet/firm_termsheet.py}{86}{103}{Estimation de la volatilité et indice cible.}\label{lst:dv:the-structured-products-business:vt}
\item \textbf{Exécuter} \texttt{dv\_structured.named\_result()}, \texttt{capped\_note()}, \texttt{reverse\_convertible()},
\texttt{vol\_target\_options()}, \texttt{decrement\_example()} et \texttt{fig\_structured.py}.
\end{tutsteps}
\textbf{Pour aller plus loin.} Abaisser la protection à 95~\% et trouver le \omterm{def:dv:the-structured-products-business:participation}{taux de participation} pour un spread de 1~\%~; remplacer le call à la monnaie par le
call à prix moyen du chapitre 16~; prélever sur l'indice cible des frais annuels de 0.5~\% et comparer le prix du call.
\end{tutorial}

\section{Construction~: le décomposeur de termes}\label{bld:dv:the-structured-products-business:termsheet}

\begin{build}
\bfield{Rôle} Le moteur de fiches de termes de la firme miniature pour son desk de \omterm{def:dv:the-structured-products-business:product}{produits structurés}~: des notes décomposées en jambes obligataires et optionnelles, les solveurs qui
traduisent la demande d'un client en conditions, et les calculateurs d'indices de stratégie.
\bfield{Interface} \texttt{Leg(kind, quantity, strike)}, \texttt{Note(maturity, legs)}~; \texttt{zero\_coupon(t, r, spread)}~;
\texttt{value(note, r, spread, q, vol\_of\_k)}~; \texttt{protected\_note(t, participation, protection, cap)}~;
\texttt{max\_participation(\ldots, margin, protection, cap)}~; \texttt{spread\_for\_participation(target, \ldots)}~;
\texttt{reverse\_convertible\_coupon}~; \texttt{ewma\_vol}, \texttt{vol\_target\_index(returns, target, rate, lam, max\_lev, dt, fee)}~;
\texttt{decrement\_index(total\_return, decrement, dt, points, base)}.
\bfield{Règles} La jambe obligataire est toujours actualisée sur la courbe de financement de l'émetteur~; les options le sont au smile du desk~; la \omterm{def:dv:the-structured-products-business:margin}{marge de structuration} est un intrant explicite,
jamais un résidu laissé dans le prix de l'option~; les indices de stratégie n'utilisent que l'information disponible à chaque rééquilibrage.
\bfield{Tests d'acceptation} \texttt{code/firm/termsheet/tests/}~: l'identité budgétaire tient et s'inverse~; un émetteur plus faible offre
un \omterm{def:dv:the-structured-products-business:participation}{taux de participation} plus élevé~; un cap achète davantage de participation~; le coupon du \omterm{def:dv:the-structured-products-business:rc}{reverse convertible} croît avec le spread et la volatilité~; un indice cible sur une
volatilité constante réalise sa cible~; un \omterm{def:dv:the-structured-products-business:decrement}{indice à décrément} sur un indice de rendement total déterministe a le forward attendu, en pourcentage et
en points.
\bfield{Pour aller plus loin} Notes sur le pricer d'\omterm{def:dv:autocallables:autocallable}{autocallable} du chapitre 18~; tarification sur le marché secondaire avec le spread courant de l'émetteur~; les
scénarios de performance du document d'informations clés.
\end{build}

\begin{omsources}
\item Règlement (UE) n°~1286/2014 du Parlement européen et du Conseil du 26 novembre 2014 relatif aux documents d'informations clés pour
les produits d'investissement packagés de détail et fondés sur l'assurance (PRIIPs), \emph{Journal officiel de l'Union européenne} L~352/1.
\item FINRA, Regulatory Notice 12-03, «~Heightened Supervision of Complex Products~» (janvier 2012).
\end{omsources}

\section{Exercices}

\begin{exercise}[$\star$]\label{exo:dv:the-structured-products-business:1}
Écrire une note protégée à 100~\% sur deux ans, avec un \omterm{def:dv:the-structured-products-business:participation}{taux de participation} de 60~\%, comme une obligation plus des options. Quelles jambes l'investisseur détient-il~?
\end{exercise}

\begin{exercise}[$\star$]\label{exo:dv:the-structured-products-business:2}
Avec un taux de 3~\%, que vaut l'obligation zéro-coupon à deux ans pour des spreads de 0 et de 2.4~\%~? Quel budget d'option le spread ajoute-t-il~?
\end{exercise}

\begin{exercise}[$\star$]\label{exo:dv:the-structured-products-business:3}
Pourquoi une note offrant un \omterm{def:dv:the-structured-products-business:participation}{taux de participation} de 70~\% n'est-elle pas nécessairement meilleure qu'une note en offrant 50~\%~?
\end{exercise}

\begin{exercise}[$\star\star$]\label{exo:dv:the-structured-products-business:4}
En utilisant les chiffres du chapitre (call 12.49, \omterm{def:dv:the-structured-products-business:margin}{marge de structuration} 1.5), quel \omterm{def:dv:the-structured-products-business:participation}{taux de participation} offre une note protégée à 95~\% pour un spread nul~?
\end{exercise}

\begin{exercise}[$\star\star$]\label{exo:dv:the-structured-products-business:5}
Expliquer pourquoi les options sur un \omterm{def:dv:the-structured-products-business:vt}{indice à volatilité cible} sont bon marché, et ce que l'investisseur abandonne.
\end{exercise}

\begin{exercise}[$\star\star$]\label{exo:dv:the-structured-products-business:6}
Un \omterm{def:dv:the-structured-products-business:decrement}{indice à décrément} de 5~\% est construit sur un indice de rendement total dont le sous-jacent verse 3~\% de dividendes. De combien retarde-t-il chaque
année sur l'indice de prix, et qui supporte l'écart si les dividendes sont ramenés à 1~\%~?
\end{exercise}

\begin{exercise}[$\star\star\star$]\label{exo:dv:the-structured-products-business:7}
\emph{Codage.} Pricer le \omterm{def:dv:the-structured-products-business:rc}{reverse convertible} de l'\cref{ex:dv:the-structured-products-business:rc} avec des prix d'exercice de 80~\%, 90~\% et 100~\% et
donner les coupons.
\end{exercise}

\begin{exercise}[$\star\star\star$]\label{exo:dv:the-structured-products-business:8}
\emph{Repérer l'erreur.} «~Nos notes offrent la meilleure participation du marché, donc notre tarification des options doit être la plus compétitive.~»
\end{exercise}

\section{Problème~: Participation à 70~\%}

\begin{problem}[{Problème du week-end --- ce qui finance la participation}]\label{pb:dv:the-structured-products-business:1}
Un distributeur demande une note de deux ans, protégée à 100~\%, indexée sur un indice avec une participation de 70~\%. Taux 3~\%, rendement du dividende 1.5~\%, surface du chapitre~9,
une \omterm{def:dv:the-structured-products-business:margin}{marge de structuration} de 1.5.

\medskip\noindent\textbf{Partie I --- La jambe optionnelle.}
\begin{enumerate}
\item Donner la volatilité à deux ans à la monnaie et le prix du call pour 100.
\item Combien coûte une participation de 70~\%~?
\item Pourquoi le call est-il évalué à la volatilité à la monnaie du smile et non à celle d'un autre prix d'exercice~?
\item Quel effet aurait sur le prix du call une moyenne du niveau final sur les six derniers mois~?
\item Que ferait un cap~?
\end{enumerate}

\medskip\noindent\textbf{Partie II --- La jambe obligataire.}
\begin{enumerate}[resume]
\item Donner la valeur de l'obligation zéro-coupon à spread nul.
\item Donner la participation disponible pour des spreads de 0, 1~\% et 2~\%.
\item Pour quel spread une participation de 70~\% est-elle possible~? Combien coûte alors l'obligation~?
\item Pourquoi le spread de l'émetteur entre-t-il dans le prix de la note~?
\item Que supporte l'investisseur en contrepartie de la participation plus élevée~?
\end{enumerate}

\medskip\noindent\textbf{Partie III --- Solutions de rechange.}
\begin{enumerate}[resume]
\item Avec un spread de 1~\%, quel cap permet une participation de 100~\%~?
\item Que permet d'acheter le même budget sous forme de \omterm{def:dv:the-structured-products-business:rc}{reverse convertible} à un an de prix d'exercice 90~\%~?
\item Quelle participation autoriserait un \omterm{def:dv:the-structured-products-business:vt}{indice à volatilité cible} de 10~\%, en prenant comme repère le rapport des volatilités implicites à un an~?
\item Pourquoi le distributeur pourrait-il préférer un \omterm{def:dv:the-structured-products-business:ssi}{indice de stratégie systématique}, et pourquoi l'investisseur pourrait-il ne pas le souhaiter~?
\item Comment le document d'informations clés encadre-t-il la présentation~?
\end{enumerate}

\medskip\noindent\textbf{Partie IV --- Jugement.}
\begin{enumerate}[resume]
\item Une participation de 70~\% à 240 points de base est-elle une bonne affaire pour le client~?
\item Que faut-il dire au distributeur au sujet d'un émetteur qui offre 70~\% alors que d'autres en offrent 50~\%~?
\item Comment un élargissement du spread de l'émetteur après l'émission affecte-t-il le prix secondaire de la note~?
\item Énoncer le \emph{résultat nommé}~: le spread de financement de l'émetteur pour lequel une note de deux ans protégée à 100~\% peut offrir une participation de 70~\% après une
\omterm{def:dv:the-structured-products-business:margin}{marge de structuration} de 1.5~\%.
\item En une phrase~: que mesure la participation d'une \omterm{def:dv:the-structured-products-business:protected}{note à capital protégé}~?
\end{enumerate}
\end{problem}

\section{Questions d'entretien}

\begin{interviewq}[$\star$ \iqroles{trader, banque}]\label{iq:dv:the-structured-products-business:1}
Comment évalueriez-vous une \omterm{def:dv:the-structured-products-business:protected}{note à capital protégé}~? Qu'est-ce qui détermine le \omterm{def:dv:the-structured-products-business:participation}{taux de participation}~?
\end{interviewq}

\begin{interviewq}[$\star$ \iqroles{banque}]\label{iq:dv:the-structured-products-business:2}
Pourquoi une \omterm{def:dv:the-structured-products-business:note}{note structurée} est-elle un instrument de financement pour la banque émettrice~?
\end{interviewq}

\begin{interviewq}[$\star\star$ \iqroles{trader}]\label{iq:dv:the-structured-products-business:3}
Qu'est-ce qu'un \omterm{def:dv:the-structured-products-business:rc}{reverse convertible}, et quelle est la position vendeuse de l'investisseur~?
\end{interviewq}

\begin{interviewq}[$\star\star$ \iqroles{chercheur}]\label{iq:dv:the-structured-products-business:4}
Pourquoi les options sur des \omterm{def:dv:the-structured-products-business:vt}{indices à volatilité cible} sont-elles moins chères que les options sur leur sous-jacent~? Quel risque subsiste pour l'émetteur~?
\end{interviewq}

\begin{interviewq}[$\star\star$ \iqroles{risque}]\label{iq:dv:the-structured-products-business:5}
Le spread de crédit de l'émetteur s'élargit de 100 points de base. Que deviennent ses notes en circulation et ses comptes~?
\end{interviewq}

\begin{interviewq}[$\star\star\star$ \iqroles{trader, banque}]\label{iq:dv:the-structured-products-business:6}
Un distributeur souhaite le coupon le plus élevé possible sur une note à un an. Passez en revue les choix qui le font monter et ce que chacun coûte à l'investisseur.
\end{interviewq}
